\color{unimain}{\subsection{Технологические защиты}}
\color{black}

\begin{enumerate}[label=\arabic{section}.\arabic{subsection}.\arabic*,  align=left, labelsep=4pt, leftmargin=0pt, itemindent=57pt]

\item
Технологические защиты от повышения давления, температуры масла и обмоток предназначены для предотвращения повреждения оборудования вследствие нарушения режима его работы, обеспечивая надёжную и безопасную работу трансформатора.   

\item
Алгоритм работы логики технологических защит трансформатора приведен на рисунке \ref{pict:TZT}, уставки функции приведены в таблице \ref{app_set:tz}.

\begin{figure}[H]
\centering
\includegraphics[width=0.9\textwidth,height=0.9\textheight,keepaspectratio]{\fbpath/07.01 РЗТ/011.1061 ТЗ Т/img/TZT.pdf}
\captionof{figure}{Функциональная схема логики ТЗ Т}\label{pict:TZT}
\end{figure}

\item
Предусмотрена возможность приема сигналов срабатывания сигнального контакта ДТм (ДТо),
отключающего контакта ДТм (ДТо) и низкой изоляции цепи отключающего контакта ДТм (ДТо), по разным
цепям от двух источников ПДС1 и ПДС2. Для приема сигналов от ПДС1 предусмотрены входные сигналы \textit{«Ср.
сигн контакта ДТм (ДТо) Т 1 цепь»}, \textit{«Ср. авар контакта ДТм (ДТо) Т 1 цепь»} и \textit{«Низкая изол. ДТм (ДТо) Т
авар 1 цепь»}, а для приема сигналов от ПДС2 предусмотрены входные сигналы \textit{«Ср. сигн контакта ДТм (ДТо)
Т 2 цепь»}, \textit{«Ср. авар контакта ДТм (ДТо) Т 2 цепь»} и \textit{«Низкая изол. ДТм (ДТо) Т авар 2 цепь»}. Если сигналы
срабатывания ДТм (ДТо) заводятся непосредственно со шкафа ШОМ силового трансформатора
конфигурируются только входные сигналы \textit{«Ср. сигн контакта ДТм (ДТо) Т 1 цепь»}, \textit{«Ср. авар контакта ДТм
(ДТо) Т 1 цепь»} и \textit{«Низкая изол. ДТм (ДТо) Т авар 1 цепь»}, а действие ДТм (ДТо) по второй цепи сигнального и
аварийного контакта выводиться из работы с помощью программных переключателей \textbf{«ДТм (ДТо) Т сигн 2
цепь»} и \textbf{«ДТм (ДТо) Т авар 2 цепь»}.

\item
Ввод в работу ДТм (ДТо) осуществляется с помощью программных переключателей \textbf{«ДТм (ДТо) Т сигн
1 цепь»}, \textbf{«ДТм (ДТо) сигн Т 2 цепь»}, \textbf{«ДТм (ДТо) Т авар 1 цепь»} и \textbf{«ДТм (ДТо) Т авра 2 цепь»}. Предусмотрена
возможность действия срабатывания аварийного контакта ДТм (ДТо) на отключение трансформатора с
помощью программных переключателей \textbf{«Откл от ДТм (ДТо) Т авар»}.

\item
Оперативный ввод/вывод  ДТм (ДТо) сигнальной и аварийной ступени производится с помощью БУ
(блока управления функциями) \textit{«Вывод ТЗ Т»}.

\item
Непрерывный контроль изоляции цепи аварийного контакта ДТм (ДТо) осуществляется с помощью реле
контроля изоляции (РКИ). В случае снижения сопротивления изоляции РКИ подает сигнал \textit{«Низкая изол. ДТм
(ДТо) Т авар 1 цепь»} (\textit{«Низкая изол. ДТм (ДТо) авар 2 цепь»}) для блокировки срабатывания ДТм (ДТо) от
аварийного контакта. Время срабатывания блокировки от сигналов РКИ задается с помощью программного
переключателя \textbf{«Тбл ТЗ Т»}. Ввод блокировки ДТм (ДТо) от РКИ осуществляется с помощью программного
переключателя \textbf{«Блок ТЗ Т»}. Для сброса блокировки от РКИ, используется команда управления \textit{«Возврат ТЗ
Т»}.


\end{enumerate} 